% =========================================================
% Style Pandoc PDF - Livrets jury QualiCheck (RNCP37827)
% Basé sur ~/.config/pandoc/styles/conception.tex (même identité visuelle)
% David LEGRAND
% =========================================================

% ---------------------------------------------------------
% Langue, typographie, micro-ajustements
% ---------------------------------------------------------
\usepackage{polyglossia}
\setmainlanguage{french}

\usepackage{microtype}
\usepackage{xcolor}
\usepackage{graphicx}
\usepackage{float}

% ---------------------------------------------------------
% Code inline (`chemin/de/fichier.md`) : autorise la coupure
% en fin de ligne. \texttt seul ne le permet pas (pas d'espaces,
% pas d'hyphénation en typo machine à écrire) — un chemin long
% déborde sinon dans la marge au lieu de passer à la ligne.
% ---------------------------------------------------------
\usepackage{seqsplit}
\let\oldtexttt\texttt
\renewcommand{\texttt}[1]{{\ttfamily\seqsplit{#1}}}

% ---------------------------------------------------------
% Couleurs sobres (identiques à conception.tex)
% ---------------------------------------------------------
\definecolor{maincolor}{HTML}{2F5597}
\definecolor{secondarycolor}{HTML}{444444}
\definecolor{lightgray}{HTML}{F4F4F4}
\definecolor{mediumgray}{HTML}{DDDDDD}
\definecolor{darkgray}{HTML}{333333}
\definecolor{codebg}{HTML}{F7F7F7}

% ---------------------------------------------------------
% Mise en page générale
% ---------------------------------------------------------
\usepackage{geometry}
\geometry{
  margin=2cm
}

\setlength{\parindent}{0pt}
\setlength{\parskip}{0.6em}
\linespread{1.08}

\clubpenalty=10000
\widowpenalty=10000

% ---------------------------------------------------------
% Titres
% ---------------------------------------------------------
\usepackage{titlesec}

\titleformat{\section}
  {\Large\bfseries\color{maincolor}}
  {\thesection}
  {0.8em}
  {}

\titleformat{\subsection}
  {\large\bfseries\color{darkgray}}
  {\thesubsection}
  {0.8em}
  {}

\titleformat{\subsubsection}
  {\normalsize\bfseries\color{secondarycolor}}
  {\thesubsubsection}
  {0.8em}
  {}

\titlespacing*{\section}{0pt}{1.4em}{0.6em}
\titlespacing*{\subsection}{0pt}{1.1em}{0.4em}
\titlespacing*{\subsubsection}{0pt}{0.9em}{0.3em}

% ---------------------------------------------------------
% Listes
% ---------------------------------------------------------
\usepackage{enumitem}

\setlist{
  itemsep=0.55em,
  topsep=0.5em,
  parsep=0.2em,
  partopsep=0pt,
  leftmargin=1.5em
}

\setlist[itemize]{label=--, itemsep=0.55em}
\setlist[enumerate]{itemsep=0.55em}
\setlist[itemize,2]{itemsep=0.35em, topsep=0.3em}
\setlist[enumerate,2]{itemsep=0.35em, topsep=0.3em}

% ---------------------------------------------------------
% Tableaux
% ---------------------------------------------------------
\usepackage{array}
\usepackage{longtable}
\usepackage{booktabs}
\usepackage{colortbl}

\renewcommand{\arraystretch}{1.25}

% ---------------------------------------------------------
% Citations Markdown
% ---------------------------------------------------------
\usepackage{framed}

\renewenvironment{quote}
  {
    \def\FrameCommand{
      {\color{maincolor}\vrule width 3pt}
      \hspace{0.8em}
    }
    \MakeFramed{\advance\hsize-\width \FrameRestore}
    \small
    \itshape
  }
  {
    \endMakeFramed
  }

% ---------------------------------------------------------
% Code inline et blocs de code
% ---------------------------------------------------------
\usepackage{fancyvrb}
\usepackage{fvextra}

\DefineVerbatimEnvironment{Highlighting}{Verbatim}{
  breaklines,
  breakanywhere,
  commandchars=\\\{\},
  fontsize=\small
}

\usepackage{mdframed}

\BeforeBeginEnvironment{verbatim}{%
  \begin{mdframed}[
    backgroundcolor=codebg,
    linecolor=mediumgray,
    linewidth=0.5pt,
    roundcorner=4pt,
    innertopmargin=0.6em,
    innerbottommargin=0.6em,
    innerleftmargin=0.8em,
    innerrightmargin=0.8em
  ]
}

\AfterEndEnvironment{verbatim}{%
  \end{mdframed}
}

% ---------------------------------------------------------
% En-têtes et pieds de page
% ---------------------------------------------------------
\usepackage{fancyhdr}
\pagestyle{fancy}

\fancyhf{}
\fancyhead[L]{QualiCheck, Livret jury}
\fancyhead[R]{\thepage}
\fancyfoot[L]{David LEGRAND}
\fancyfoot[R]{\today}

\renewcommand{\headrulewidth}{0.4pt}
\renewcommand{\footrulewidth}{0.4pt}

% ---------------------------------------------------------
% Liens
% ---------------------------------------------------------
\usepackage{hyperref}

\hypersetup{
  colorlinks=true,
  linkcolor=maincolor,
  urlcolor=maincolor,
  citecolor=maincolor,
  pdfborder={0 0 0}
}

% ---------------------------------------------------------
% Images
% ---------------------------------------------------------
\setkeys{Gin}{width=\linewidth, keepaspectratio}

% ---------------------------------------------------------
% Table des matières
% ---------------------------------------------------------
\setcounter{tocdepth}{3}

% =========================================================
% Page de garde — commune aux 6 livrets (Livret 0 + E1 à E5)
%
% Remplace \maketitle : construite depuis le Front Matter Pandoc
% (title, subtitle, author, date). \@title porte déjà le sous-titre
% (mécanisme \subtitle du template par défaut de Pandoc, qui l'ajoute
% à \@title avant que \maketitle ne s'exécute).
%
% \devise : optionnelle, vide par défaut — seul le livret concerné la
% renseigne, dans le Front Matter du .md (exécuté en préambule, donc avant
% \maketitle — un bloc LaTeX brut dans le corps du document s'exécuterait
% trop tard, après que la page de garde est déjà composée) :
%   header-includes: |
%     \def\devise{Garbage in, garbage out}
% =========================================================
\usepackage{etoolbox}
\def\devise{}

\makeatletter
\renewcommand{\maketitle}{%
  \begin{titlepage}
    \thispagestyle{empty}
    \centering
    \vspace*{3cm}

    {\color{maincolor}\rule{\linewidth}{1.2pt}}

    \vspace{1.5cm}
    {\Huge\bfseries\color{maincolor} \@title \par}
    \vspace{2cm}

    {\color{maincolor}\rule{\linewidth}{1.2pt}}

    \ifdefempty{\devise}{}{%
      \vspace{1em}
      {\normalsize\itshape\color{secondarycolor} \devise\par}
    }

    \vfill

    {\Large \@author \par}
    \vspace{0.5em}
    {\large \@date \par}
    \vspace{2cm}

    {\normalsize\color{secondarycolor}
      Certification RNCP37827, Développeur en intelligence artificielle (saveur agentique)\par
      Projet QualiCheck, Simplon.co\par}

    \vspace{2cm}
  \end{titlepage}
  \newpage
}
\makeatother
